% Options for packages loaded elsewhere
% Options for packages loaded elsewhere
\PassOptionsToPackage{unicode}{hyperref}
\PassOptionsToPackage{hyphens}{url}
\PassOptionsToPackage{dvipsnames,svgnames,x11names}{xcolor}
%
\documentclass[
  letterpaper,
  DIV=11,
  numbers=noendperiod]{scrartcl}
\usepackage{xcolor}
\usepackage{amsmath,amssymb}
\setcounter{secnumdepth}{-\maxdimen} % remove section numbering
\usepackage{iftex}
\ifPDFTeX
  \usepackage[T1]{fontenc}
  \usepackage[utf8]{inputenc}
  \usepackage{textcomp} % provide euro and other symbols
\else % if luatex or xetex
  \usepackage{unicode-math} % this also loads fontspec
  \defaultfontfeatures{Scale=MatchLowercase}
  \defaultfontfeatures[\rmfamily]{Ligatures=TeX,Scale=1}
\fi
\usepackage{lmodern}
\ifPDFTeX\else
  % xetex/luatex font selection
\fi
% Use upquote if available, for straight quotes in verbatim environments
\IfFileExists{upquote.sty}{\usepackage{upquote}}{}
\IfFileExists{microtype.sty}{% use microtype if available
  \usepackage[]{microtype}
  \UseMicrotypeSet[protrusion]{basicmath} % disable protrusion for tt fonts
}{}
\makeatletter
\@ifundefined{KOMAClassName}{% if non-KOMA class
  \IfFileExists{parskip.sty}{%
    \usepackage{parskip}
  }{% else
    \setlength{\parindent}{0pt}
    \setlength{\parskip}{6pt plus 2pt minus 1pt}}
}{% if KOMA class
  \KOMAoptions{parskip=half}}
\makeatother
% Make \paragraph and \subparagraph free-standing
\makeatletter
\ifx\paragraph\undefined\else
  \let\oldparagraph\paragraph
  \renewcommand{\paragraph}{
    \@ifstar
      \xxxParagraphStar
      \xxxParagraphNoStar
  }
  \newcommand{\xxxParagraphStar}[1]{\oldparagraph*{#1}\mbox{}}
  \newcommand{\xxxParagraphNoStar}[1]{\oldparagraph{#1}\mbox{}}
\fi
\ifx\subparagraph\undefined\else
  \let\oldsubparagraph\subparagraph
  \renewcommand{\subparagraph}{
    \@ifstar
      \xxxSubParagraphStar
      \xxxSubParagraphNoStar
  }
  \newcommand{\xxxSubParagraphStar}[1]{\oldsubparagraph*{#1}\mbox{}}
  \newcommand{\xxxSubParagraphNoStar}[1]{\oldsubparagraph{#1}\mbox{}}
\fi
\makeatother


\usepackage{longtable,booktabs,array}
\usepackage{calc} % for calculating minipage widths
% Correct order of tables after \paragraph or \subparagraph
\usepackage{etoolbox}
\makeatletter
\patchcmd\longtable{\par}{\if@noskipsec\mbox{}\fi\par}{}{}
\makeatother
% Allow footnotes in longtable head/foot
\IfFileExists{footnotehyper.sty}{\usepackage{footnotehyper}}{\usepackage{footnote}}
\makesavenoteenv{longtable}
\usepackage{graphicx}
\makeatletter
\newsavebox\pandoc@box
\newcommand*\pandocbounded[1]{% scales image to fit in text height/width
  \sbox\pandoc@box{#1}%
  \Gscale@div\@tempa{\textheight}{\dimexpr\ht\pandoc@box+\dp\pandoc@box\relax}%
  \Gscale@div\@tempb{\linewidth}{\wd\pandoc@box}%
  \ifdim\@tempb\p@<\@tempa\p@\let\@tempa\@tempb\fi% select the smaller of both
  \ifdim\@tempa\p@<\p@\scalebox{\@tempa}{\usebox\pandoc@box}%
  \else\usebox{\pandoc@box}%
  \fi%
}
% Set default figure placement to htbp
\def\fps@figure{htbp}
\makeatother


% definitions for citeproc citations
\NewDocumentCommand\citeproctext{}{}
\NewDocumentCommand\citeproc{mm}{%
  \begingroup\def\citeproctext{#2}\cite{#1}\endgroup}
\makeatletter
 % allow citations to break across lines
 \let\@cite@ofmt\@firstofone
 % avoid brackets around text for \cite:
 \def\@biblabel#1{}
 \def\@cite#1#2{{#1\if@tempswa , #2\fi}}
\makeatother
\newlength{\cslhangindent}
\setlength{\cslhangindent}{1.5em}
\newlength{\csllabelwidth}
\setlength{\csllabelwidth}{3em}
\newenvironment{CSLReferences}[2] % #1 hanging-indent, #2 entry-spacing
 {\begin{list}{}{%
  \setlength{\itemindent}{0pt}
  \setlength{\leftmargin}{0pt}
  \setlength{\parsep}{0pt}
  % turn on hanging indent if param 1 is 1
  \ifodd #1
   \setlength{\leftmargin}{\cslhangindent}
   \setlength{\itemindent}{-1\cslhangindent}
  \fi
  % set entry spacing
  \setlength{\itemsep}{#2\baselineskip}}}
 {\end{list}}
\usepackage{calc}
\newcommand{\CSLBlock}[1]{\hfill\break\parbox[t]{\linewidth}{\strut\ignorespaces#1\strut}}
\newcommand{\CSLLeftMargin}[1]{\parbox[t]{\csllabelwidth}{\strut#1\strut}}
\newcommand{\CSLRightInline}[1]{\parbox[t]{\linewidth - \csllabelwidth}{\strut#1\strut}}
\newcommand{\CSLIndent}[1]{\hspace{\cslhangindent}#1}



\setlength{\emergencystretch}{3em} % prevent overfull lines

\providecommand{\tightlist}{%
  \setlength{\itemsep}{0pt}\setlength{\parskip}{0pt}}



 


\KOMAoption{captions}{tableheading}
\makeatletter
\@ifpackageloaded{caption}{}{\usepackage{caption}}
\AtBeginDocument{%
\ifdefined\contentsname
  \renewcommand*\contentsname{Table of contents}
\else
  \newcommand\contentsname{Table of contents}
\fi
\ifdefined\listfigurename
  \renewcommand*\listfigurename{List of Figures}
\else
  \newcommand\listfigurename{List of Figures}
\fi
\ifdefined\listtablename
  \renewcommand*\listtablename{List of Tables}
\else
  \newcommand\listtablename{List of Tables}
\fi
\ifdefined\figurename
  \renewcommand*\figurename{Figure}
\else
  \newcommand\figurename{Figure}
\fi
\ifdefined\tablename
  \renewcommand*\tablename{Table}
\else
  \newcommand\tablename{Table}
\fi
}
\@ifpackageloaded{float}{}{\usepackage{float}}
\floatstyle{ruled}
\@ifundefined{c@chapter}{\newfloat{codelisting}{h}{lop}}{\newfloat{codelisting}{h}{lop}[chapter]}
\floatname{codelisting}{Listing}
\newcommand*\listoflistings{\listof{codelisting}{List of Listings}}
\makeatother
\makeatletter
\makeatother
\makeatletter
\@ifpackageloaded{caption}{}{\usepackage{caption}}
\@ifpackageloaded{subcaption}{}{\usepackage{subcaption}}
\makeatother
\usepackage{bookmark}
\IfFileExists{xurl.sty}{\usepackage{xurl}}{} % add URL line breaks if available
\urlstyle{same}
\hypersetup{
  pdftitle={Childcare Costs \& Female Labor Force Participation},
  colorlinks=true,
  linkcolor={blue},
  filecolor={Maroon},
  citecolor={Blue},
  urlcolor={Blue},
  pdfcreator={LaTeX via pandoc}}


\title{Childcare Costs \& Female Labor Force Participation}
\usepackage{etoolbox}
\makeatletter
\providecommand{\subtitle}[1]{% add subtitle to \maketitle
  \apptocmd{\@title}{\par {\large #1 \par}}{}{}
}
\makeatother
\subtitle{DSAN 5650: Causal Inference for Computational Social Science}
\author{Emily Mitchum}
\date{2026-08-10}
\begin{document}
\maketitle
\begin{abstract}
A substantial body of research has found that reductions in childcare
costs are positively associated with increased maternal labor-force
participation. However, when analyzing the elasticity of employment to
childcare prices, researchers find heterogeneity across various
geographies and time-periods. This may be attributable to differences in
educational attainment, labor force attachment, and other factors
(Morrissey 2016). This research paper uses U.S. county-level data to
examine whether the association between childcare affordability and
maternal labor-force participation among women with children under age
six varies across local occupational contexts. I present a Bayesian
logit-normal regression model that estimates maternal labor force
participation as a function of childcare affordability, the county share
of workers in ``greedy'' occupations, educational attainment, geographic
rootedness, and an interaction between childcare burden and occupational
structure. Aligning with existing research, I find that higher childcare
burden is associated with lower maternal labor force participation.
However, I also find that this negative association becomes weaker as a
county's concentration of women in greedy occupations increases.
Educational attainment and geographic rootedness are also moderately
associated with maternal labor force participation, while
greedy-occupation-share is weakly associated.
\end{abstract}


\subsection{Introduction}\label{sec-1}

Earlier this year, I examined the major factors that help explain U.S.
maternal labor force participation trends and persistent wage
differentials between men and women. Existing research consistently
identifies motherhood as an important contributor to both reduced
labor-force participation and long-term earnings penalties.

Childcare costs, which have significantly outpaced overall inflation,
play a significant role in women's ability or decision to remain in the
workforce. This trend has driven an overall focus within the research
community on the relationship between childcare affordability and
maternal labor force participation. Research has appropriately paid
particular attention to low-income women, for whom childcare expenses
represent a large share of potential earnings, and whose labor market
attachment is the most impacted by childcare costs
(https://www.jstor.org/stable/41238095).

Several studies have concluded that reducing childcare costs would
increase maternal employment, especially among low-income women. The
majority of U.S. childcare assistance programs are also primarily
available for lower-income households and reach only a minority of
eligible families. Less is known about this association amongst women in
moderate or high income brackets, who remain burdened by high childcare
costs. For these women, childcare constraints may impact labor force
participation differently. Increased earning potential and financial
resources may make them more likely to remain employed despite high
childcare costs, while reducing hours or seeking work flexibility to
meet childcare demands. This behavior may ultimately result in reduced
career advancement and lower wages.

This paper tests this theory, using county-level data to analyze how
women in different occupational contexts may respond differently to
childcare costs. I use the concentration of women working in ``greedy''
occupations, or high-paying occupations that disproportionately reward
long hours and constant availability, to measure local occupational
context. This is used to test whether the known association between
childcare burden and maternal labor-force participation is moderated by
occupation ``greediness''.

County-level female labor force participation only determines whether
mothers participate in the labor force, rather than how many hours they
work. Therefore, this analysis does not directly model the wage gap.
However, beginning to understand potential effect heterogeneity in this
relationship may help explain persistent gender gaps in earnings despite
comparable levels of overall labor-force participation.

\subsection{Background and Conceptual
Framework}\label{background-and-conceptual-framework}

\subsubsection{Existing Childcare Policy and Maternal
Employment}\label{existing-childcare-policy-and-maternal-employment}

Researchers generally agree that increased investment in childcare
affordability interventions results in higher maternal labor force
participation, particularly among low-income women in the first few
years after birth. For instance, Burgess, Chien and Enchautegui find
that increasing Child Care and Development Block Grant (CCDBG)
expenditures by 10-percent per child would increase employment levels
among low-income women with young children by between half to two-thirds
of a percent
(https://aspe.hhs.gov/sites/default/files/migrated\_legacy\_files//171051/EffectsCCSubsidiesMaternalLFPBrief.pdf).

Still, the majority of U.S. childcare assistance programs reach only a
subset of families and are concentrated among lower-income households.
The Child Care and Development Block Grant (CCDBG) is only available for
families who make less than 85\% of their state median income. Even
amongst those that meet requirements, federal childcare subsidies only
reach approximately 14\% of eligible families
(https://aspe.hhs.gov/reports/factsheet-estimates-child-care-eligibility-receipt-fiscal-year-2017),
and state subsidies reach only approximately 22\% of eligible families.
Tax incentives, including the Child and Dependent Care Tax Credit, also
benefit a minority of families with children due to its non-refundable
structure, steep income phase-downs, and restrictions that require two
working parents.
(https://www2.census.gov/library/working-papers/2025/adrm/ces/CES-WP-25-25.pdf)

While most childcare programs focus on low-income families, evidence
from universal early childcare programs suggests that reducing childcare
constraints can affect maternal employment across a broader range of
women. Two years after implementing a universal pre-kindergarten program
in 2008, D.C. observed a 12\% increase in the maternal labor force
participation rate. A study from the Center for American Progress found
that this increase was strongly associated with implementation of
universal preschool
(https://www.americanprogress.org/article/effects-universal-preschool-washington-d-c/).
This example is particularly important given D.C.'s occupational
profile, with an unusually large concentration of highly-educated women
in high-paying, senior positions. This example motivates the current
analysis, which seeks to understand the varying slope of the
relationship between childcare affordability and labor force
participation among counties with different female occupational
profiles.

\subsubsection{Occupational
``Greediness''}\label{occupational-greediness}

This analysis builds on existing research from economist Claudia Goldin,
who has extensively studied ``greedy occupations''. Goldin finds that
greedy occupations lead women to remain employed, even as they have
children, but earn less over time, due to childcare responsibilities
that prevent them from meeting demanding work requirements, including
long hours and travel (Goldin 2014). Women's role as primary care-givers
leads them to take career breaks, reduce hours, and trade high-stress
roles for flexible roles.

Goldin identifies finance, law, consulting, and corporate management as
common greedy occupations, characterized by time inflexibility, low work
substitutability, and non-linear pay. These occupations drive gender pay
differentials, since women are more likely to experience childcare
related work interruptions and less likely to be able to work rigid
schedules (Sobeck 2024). Using Goldin's research as a framework, I model
county-level occupational greediness by sex using U.S. Census adult
employment statistics.

\subsubsection{Educational Attainment}\label{educational-attainment}

Educational attainment is closely related to occupation context and
labor force participation. Geographies with a greater share of
individuals with advanced degrees are more likely to have higher
concentrations of individuals in advanced white-collar occupations.
Women in these occupations are more likely to have greater earning
power, which also influences their labor force attachment. I model
educational attainment as the share of individuals in each county with
at least a bachelor's degree. This is used to determine if the effect of
greedy occupation structure is partially or completely explained by
county educational composition.

\subsubsection{Geographic ``Rootedness'' and Informal Childcare
Support}\label{geographic-rootedness-and-informal-childcare-support}

Informal childcare support, through family-members, friends and
neighbors, can provide mothers with relief from exorbitant childcare
costs, as well as unmatched flexibility. However, the availability of
local support varies across households. Highly-educated and high-income
individuals are often more likely to move away from their birthplace,
where family resides, to regions with greater economic opportunities for
higher-paying jobs. While over half of U.S. adults report living within
an hour of family, a Pew Research study found that upper-income adults
and adults with a postgraduate degree were the most likely to report
living near no extended family.
https://www.pewresearch.org/short-reads/2022/05/18/more-than-half-of-americans-live-within-an-hour-of-extended-family/

Existing literature suggests that access to informal familial childcare
can influence maternal labor force participation. Compton and Pollak
find that married women with young children that live within close
proximity of mothers or mothers-in-law are predicted to have a 4 to 10
percent higher labor force participation rate when compared to women
that live farther from family
(https://www.sciencedirect.com/science/article/pii/S0094119013000314).
They theorize that family members within close proximity provide
``insurance'' to working mothers, by being available to provide care in
irregular or unanticipated circumstances. Building upon this framework,
I use geographic rootedness as a proxy for access to local familial
childcare. I model geographic ``rootedness'' as the share of adults 25
or older in each county who report living in the state they were born
in. While this does not directly estimate the proximity of family or
availability of informal childcare, it provides a county-level profile
of geographic attachment to measure potential differences in access to
informal childcare support.

\subsection{Data and Measures}\label{data-and-measures}

\subsubsection{Data Sources}\label{data-sources}

This analysis uses 2022 data from the Department of Labor's National
Database of Childcare Prices (NDCP) to retrieve each county's maternal
female labor force participation rate (FLPR) for women with children
under 6, the outcome variable, and childcare costs, which are used to
determine the childcare cost burden. The American Community Survey (ACS)
five-year estimates are used to identify the median family income, share
of bachelor's-plus individuals, female greedy occupation share, and
geographic rootedness. This county level data is merged using
FIPS-codes, resulting in a final sample of 2,653 counties. 491 counties
are excluded from the sample as a result of missing any of the modeled
variables. Exclusion is primarily a result of missing childcare data,
including complete state absences in Pennsylvania, Indiana, Vermont, and
New Mexico, due to low response rates or too few childcare providers.
This is a known limitation of the NDCP.

\subsubsection{Data Transformations}\label{data-transformations}

FLPR is first converted from a percentage to a proportion. Exact 0 and 1
FLPR values are then slightly adjusted, since the logit function is
undefined at those exact values. Greedy occupation share is represented
by the proportion of eligible employed women working in greedy
occupations. Bachelor's-plus share is calculated as the proportion of
adults age 25 and older with at least a bachelor's degree. Finally,
geographic rootedness is calculated as the proportion of adults age 25
and older who were born in their current state of residence. All
continuous predictors are standardized to have mean zero and standard
deviation prior to modeling.

\subsubsection{Data and Measures
Summary}\label{data-and-measures-summary}

The sources and representations of each variable are summarized below:
\textbar{} Variable \textbar{} Source(s) \textbar{} Representation
\textbar{} \textbar---\textbar---\textbar---\textbar{} \textbar{}
Maternal Female Labor Force Participation Rate (FLPR) \textbar{}
Department of Labor (DOL) county data \textbar{} Labor-force
participation rate among women with children under age six \textbar{}
\textbar{} Childcare cost burden \textbar{} DOL National Database of
Childcare Prices; American Community Survey (ACS) median family income
\textbar{} Annual center-based childcare cost for children ages 0--5
divided by median family income \textbar{} \textbar{} Greedy occupation
share \textbar{} ACS Occupation by Sex for the Civilian Employed
Population 16 Years and Over (Series S240) \textbar{} Share of employed
women in management, business and financial, and legal occupations
\textbar{} \textbar{} Educational attainment \textbar{} ACS Place of
Birth by Educational Attainment in the United States (Series B06009)
\textbar{} Share of adults age 25 or older with a bachelor's, graduate,
or professional degree \textbar{} \textbar{} Geographic rootedness
\textbar{} ACS Place of Birth by Educational Attainment in the United
States (Series B06009) \textbar{} Share of adults age 25 or older who
were born in their current state of residence \textbar{}

\subsection{Empirical Strategy}\label{empirical-strategy}

\subsubsection{Naive Model: FLPR as a Function of Childcare
Burden}\label{naive-model-flpr-as-a-function-of-childcare-burden}

First, I analyze a naive logit-Normal regression model of female labor
force participation as a function of only childcare burden.

This relationship can be written as:
\begin{equation}\phantomsection\label{eq-1}{
y_i = \beta_0 + \beta_Cc_i
}\end{equation}

where \(i\) indexes counties, \(y_i\) is the female labor force
participation rate, \(b_0\) represents the baseline female labor force
participation rate when childcare burden is held constant, or at an
average childcare burden across counties, and \(b_C\) represents the
strength of the relationship between childcare cost burden and female
labor force participation.

FLPR is modeled using a normal likelihood. The intercept term, or
baseline female labor force participation rate, is drawn from a normal
distribution with a mean of 0.405, or \(logit(70)\), since we assume the
rate is close to the mean of 70\% across counties.

Figure 1 shows a Bayesian probabilistic graphical model (PGM) depicting
the Naive model.

\pandocbounded{\includegraphics[keepaspectratio]{index_files/mediabag/assets/naive_model.pdf}}

\subsubsection{Adjusting for Covariates}\label{adjusting-for-covariates}

Next, I fit progressively more complex models, adding relevant
covariates and comparing coefficient estimates. All models are estimated
using Bayesian logit-Normal regression. This approach allows me to
assess robustness of the estimated relationships after adjusting for
educational attainment and geographic rootedness and whether these
variables account for part of the association attributed to greedy
occupational structure.

\subsubsection{Hypotheses}\label{hypotheses}

My predictions about the relationships between county-level childcare
cost burden, greedy occupation share, educational attainment, and
geographic rootedness yield the following hypotheses:

\begin{itemize}
\tightlist
\item
  \textbf{Hypothesis 1.} \(\beta_{C} < 0\) (\emph{hypothesized causal
  mechanism}: higher childcare cost burden is associated with lower
  county-level maternal labor-force participation).
\item
  \textbf{Hypothesis 2.} \(\beta_{CG} > 0\) (\emph{hypothesized causal
  mechanism}: the negative association between childcare burden and
  maternal labor-force participation weakens as greedy occupation share
  increases).
\item
  \textbf{Hypothesis 3.} \(\beta_{B} > 0\) (\emph{hypothesized causal
  mechanism}: inclusion of educational attainment reduces the estimated
  coefficient of greedy occupation share, indicating that educational
  composition accounts for part of the initially observed association).
\item
  \textbf{Hypothesis 4.} \(\beta_{R} > 0\) (\emph{hypothesized causal
  mechanism}: county-level geographic rootedness is associated with
  higher maternal labor-force participation, potentially because
  rootedness is associated with stronger access to informal family
  support networks).
\end{itemize}

The first hypothesis, captured by the sign of \(\beta_1\), is that
ruggedness has a direct negative effect on income. The second
hypothesis, which is core to our paper, is that in Africa only,
ruggedness has an additional positive effect on income. The third
hypothesis compares the estimated direct effects of ruggedness in the
unrestricted and restricted equations. We expect that, since \(\beta_5\)
absorbs both the direct and indirect effects of ruggedness, it will be
an upwards biased estimate of \(-\alpha\), whereas \(\beta_1\) will be a
consistent estimate of \(-\alpha\). From this it follows that
\(\beta_5\) is greater than \(\beta_1\). The following section tests
whether the data supports these three hypotheses.

\textbf{?@eq-5} illustrates the relationship between income and
ruggedness, leaving slave exports in the background. To bring slave
exports to the foreground and check to what extent they can account for
the differential effect of ruggedness in Africa, in section 5 we derive
a different reduced-form equation. Consider again our three equations
(\ref{eq-1}) to (\textbf{?@eq-3}) that give the relationships between
ruggedness, slave exports, and income. Previously, we substituted both
(\textbf{?@eq-3}) and (\textbf{?@eq-2}) into (\ref{eq-1}) to obtain a
reduced-form relationship between income and ruggedness only, given by
(\textbf{?@eq-4}). If we instead only substitute (\textbf{?@eq-2}) into
(\ref{eq-1}), we obtain a reduced-form relationship between income and
both ruggedness and slave exports:

\[
y_i = \begin{cases}
\kappa_1 - \alpha r_i + \beta \kappa_2 - \beta \gamma x_i + \zeta_i &\text{if }i\text{ is in Africa}, \\
\kappa_1 - \alpha r_i + \zeta_i &\text{otherwise}.
\end{cases}
\]

We test this relationship by estimating the following equation (note
that for all non-African countries slave exports are zero, \(x_i = 0\)):

\begin{equation}\phantomsection\label{eq-7}{
y_i = \beta_6 + \beta_7 r_i + \beta_8 r_i I_i^\text{Africa} + \beta_9 I_i^\text{Africa} + \beta_{10}x_i + \varepsilon_i
}\end{equation}

Estimating this reduced form allows us to test three additional
hypotheses that are about ruggedness and slave exports:

\begin{itemize}
\tightlist
\item
  Hypothesis 4. \(\beta_{10} < 0\) (slave exports negatively affect
  income).
\item
  Hypothesis 5. \(\beta_8 = 0\) (once slave exports are taken into
  account, the effect of ruggedness is no different in Africa).
\item
  Hypothesis 6. \(\beta_7 < 0\) (once slave exports are taken into
  account, the effect of ruggedness is negative).
\end{itemize}

Hypothesis 4 is that slave exports are negatively related to current
income. Hypothesis 5 provides a way of testing whether the slave trades
can fully account for the positive indirect effect of ruggedness in
Africa. If the ruggedness-income relationship is different for Africa
only because of the slave trades, then once we control for the effect of
the slave trades on income, there should no longer be a differential
effect of ruggedness for Africa. The last hypothesis is that once the
indirect effect of ruggedness is taken into account by controlling for
the slave trades, the direct effect of ruggedness can be correctly
identified and should be negative.

\subsection{The direct and indirect effects of ruggedness}\label{sec-4}

As a first step in our empirical analysis, we now estimate the direct
and indirect effects of ruggedness on income per person. Our baseline
estimates of equations (\textbf{?@eq-5}) and (\textbf{?@eq-6}) are given
in Table~\ref{tbl-1}.

\begin{longtable}[]{@{}
  >{\raggedright\arraybackslash}p{(\linewidth - 14\tabcolsep) * \real{0.1250}}
  >{\raggedright\arraybackslash}p{(\linewidth - 14\tabcolsep) * \real{0.1250}}
  >{\raggedright\arraybackslash}p{(\linewidth - 14\tabcolsep) * \real{0.1250}}
  >{\raggedright\arraybackslash}p{(\linewidth - 14\tabcolsep) * \real{0.1250}}
  >{\raggedright\arraybackslash}p{(\linewidth - 14\tabcolsep) * \real{0.1250}}
  >{\raggedright\arraybackslash}p{(\linewidth - 14\tabcolsep) * \real{0.1250}}
  >{\raggedright\arraybackslash}p{(\linewidth - 14\tabcolsep) * \real{0.1250}}
  >{\raggedright\arraybackslash}p{(\linewidth - 14\tabcolsep) * \real{0.1250}}@{}}

\caption{\label{tbl-1}The direct and indirect effects of ruggedness}

\tabularnewline

\toprule\noalign{}
\begin{minipage}[b]{\linewidth}\raggedright
\end{minipage} &
\multicolumn{7}{>{\raggedright\arraybackslash}p{(\linewidth - 14\tabcolsep) * \real{0.8750} + 12\tabcolsep}@{}}{%
\begin{minipage}[b]{\linewidth}\raggedright
Dependent variable: Log real gdp per person 2000
\end{minipage}} \\
\begin{minipage}[b]{\linewidth}\raggedright
\end{minipage} & \begin{minipage}[b]{\linewidth}\raggedright
(1)
\end{minipage} & \begin{minipage}[b]{\linewidth}\raggedright
(2)
\end{minipage} & \begin{minipage}[b]{\linewidth}\raggedright
(3)
\end{minipage} & \begin{minipage}[b]{\linewidth}\raggedright
(4)
\end{minipage} & \begin{minipage}[b]{\linewidth}\raggedright
(5)
\end{minipage} & \begin{minipage}[b]{\linewidth}\raggedright
(6)
\end{minipage} & \begin{minipage}[b]{\linewidth}\raggedright
(7)
\end{minipage} \\
\midrule\noalign{}
\endhead
\bottomrule\noalign{}
\endlastfoot
Ruggedness & \begin{minipage}[t]{\linewidth}\raggedright
-0.078\\
{(0.036)**}\strut
\end{minipage} & \begin{minipage}[t]{\linewidth}\raggedright
0.003\\
{(0.031)}\strut
\end{minipage} & \begin{minipage}[t]{\linewidth}\raggedright
-0.075\\
{(0.037)**}\strut
\end{minipage} & \begin{minipage}[t]{\linewidth}\raggedright
-0.094\\
{(0.036)***}\strut
\end{minipage} & \begin{minipage}[t]{\linewidth}\raggedright
-0.073\\
{(0.029)**}\strut
\end{minipage} & \begin{minipage}[t]{\linewidth}\raggedright
-0.081\\
{(0.031)***}\strut
\end{minipage} & \begin{minipage}[t]{\linewidth}\raggedright
-0.093\\
{(0.031)***}\strut
\end{minipage} \\
Observations & 170 & 170 & 170 & 170 & 170 & 170 & 170 \\

\end{longtable}

Looking first at column (1), when we estimate equation (\textbf{?@eq-5})
by regressing income per person on ruggedness while allowing for a
differential effect in African countries, we find empirical confirmation
for our hypotheses 1 and 2: ruggedness has a direct negative effect on
income per person, but within Africa, there is an additional positive
effect. The coefficient for ruggedness is negative and statistically
significant (i.e., \(\beta_1 < 0\) in equation (\textbf{?@eq-5})), while
the coefficient for ruggedness interacted with an indicator variable for
Africa is positive and statistically significant (i.e., β2
\textgreater{} 0 in equation (5)). According to the estimated
coefficients, within Africa the indirect positive effect swamps the
direct negative effect. Adding up the negative direct effect (-0.078)
and the positive indirect effect (0.156) yields a total effect of
ruggedness on income in African countries of 0.078. Hence, overall,
ruggedness has been a blessing for African countries.

In terms of magnitude, within Africa, a one standard deviation increase
in ruggedness increases log real gdp per person by .26 standard
deviations. This is the combination of a reduction of income per person
by .26 standard deviations due to the negative direct effect and an
increase by .51 standard deviations due to the positive indirect effect
for Africa. Put differently, for a country at the mean African level of
real income per person (\$1,784), a one standard deviation increase in
ruggedness relative to the African mean (an increase in average uphill
slope from 2.69\% to 5.74\%) is estimated to raise income to \$2,264,
which is an increase of \$481. This magnitude is quite substantial,
especially given that we are considering one very specific geographic
characteristic.

Next, in column (2), we repeat the estimation restricting the effect of
ruggedness on income to be the same for African and non-African
countries. Consistent with hypothesis 3 that the negative direct effect
of ruggedness is biased upwards towards zero when the positive effect of
ruggedness in Africa is not taken into account, the estimated
coefficient of ruggedness in column (2) is greater than the coefficient
in column (1) (i.e., \(\beta_5\) in equation (\textbf{?@eq-6}) is
greater than \(\beta_1\) in equation (\textbf{?@eq-5})).

\subsubsection{Robustness with respect to omitted geographical
variables}\label{robustness-with-respect-to-omitted-geographical-variables}

When interpreting our core results regarding the relationship between
ruggedness and current economic outcomes, a possible source of concern
is that they may be driven, at least in part, by other geographical
features correlated with ruggedness. Our primary interest is in the
differential effect of ruggedness for Africa. In order for an omitted
variable to bias this estimate, it must be the case that, either the
relationship between income and the omitted factor is different within
Africa and outside of Africa, or the relationship between the omitted
factor and ruggedness is different within and outside of Africa.

A possible curse of mineral resources (Sachs and Warner, 2001, Mehlum,
Moene, and Torvik, 2006b) could potentially affect our results for
precisely this reason. If diamond deposits are correlated with
ruggedness, and diamond production increases income outside of Africa,
but decreases income within Africa because of poor institutions, then
this could potentially bias the estimated differential effect of
ruggedness.\footnote{See Mehlum et al.~(2006b) and Robinson, Torvik, and
  Verdier (2006) for theory and empirical evidence supporting such a
  differential effect of resource endowments.} To deal with potentially
omitted differential effects such as this, we must add to our baseline
specification of column (1) both the control variable (diamond
production 1994--2000 normalized by land area in this case) and an
interaction of the control variable with our Africa indicator variable
to allow the effect of the control variable to differ for Africa. Column
(3) shows that the effect of diamonds is positive and statistically
significant in general, but that for African countries there is a
differential negative and statistically significant effect that nearly
wipes out the general positive effect. Nevertheless, the inclusion of
this control variable and its interaction with the Africa indicator
variable does not alter our results regarding the relationship between
ruggedness and current economic outcomes.\footnote{10We have also tried
  including measures of other mineral resources, oil reserves and gold
  in particular, again together with an African interaction term but
  found no significant effects.}

A related argument can be made about tropical diseases. If rugged areas
are less prone to tropical diseases within Africa but not in the rest of
the world, then this could potentially bias the estimated differential
effect of ruggedness. To check this possibility, in column (4) we add to
our baseline specification of column (1) a variable measuring the
percentage of each country that has any of the four tropical climates in
the Köppen-Geiger climate classification, as well as an interaction of
this variable with the Africa indicator variable. We see that there is a
statistically significant negative relationship between tropical climate
and income, but that the effect is no different for African countries.
Our core results are, once again, unchanged.\footnote{If instead of
  looking at tropical diseases in general, we focus on malaria in
  particular by including an index of the stability of malaria
  transmission from Kiszewski, Mellinger, Spielman, Malaney, Sachs, and
  Sachs (2004) and the corresponding African interaction, our core
  results are also unchanged. The same is true if we include distance to
  the equator and the corresponding African interaction.}

Given the arguments of Mellinger et al.~(2000) and others that coastal
access and soil quality are fundamental determinants of world income
differences, in columns (5) and (6) we control respectively for the
average distance to the nearest ice-free coast for each country (in
thousands of kilometres) and for the percentage of each country's land
that is desert. As before, we also include an interaction of each of
these variables with the African indicator variable. Our results remain
robust. Finally, in column (7) we include all these geographic controls
and their corresponding interaction terms. We find again that our
baseline results from column (1) are robust to controlling for other
geographic characteristics that could have a differential effect in
Africa.

\subsubsection{Robustness with respect to influential
observations}\label{robustness-with-respect-to-influential-observations}

Next, we check whether the results from Table~\ref{tbl-1} are driven by
some particularly influential outliers. Figure 2 shows a scatter plot of
income per person against ruggedness for African countries (top panel)
and non-African countries (bottom panel). In these plots of the raw
data, one observes a positive relationship for African countries and a
negative relationship for non-African countries. However, a number of
observations appear as clear outliers in terms of their ruggedness,
particularly, Lesotho and the Seychelles in Africa, and Switzerland,
Nepal, Tajikistan, and Kyrgyzstan outside of Africa. Our first
sensitivity check estimates our baseline specification after dropping
these six observations, without and with our full set of control
variables. Results are presented in columns (1) and (2) of table
2.\footnote{In the scatter plot, one can also observe several small
  countries that have either unusually high or unusually low ruggedness.
  These observations are not particularly influential for our results,
  which remain qualitatively unchanged if we drop the 20 countries in
  our regressions with an area below 5,000 square kilometres (which is
  slightly above the area of Seychelles) or if we estimate a weighted
  least squares regression using countries' areas as weights. A related
  concern is that our results may be driven by ``atypical'' African
  countries, such as island countries or North African countries. Our
  results are also robust to omitting these ``atypical'' African
  countries from the sample.}

We next adopt a more systematic approach to deal with influential
observations, and remove influential observations using each
observation's \texttt{DFBETA}, which is a measure of the difference in
the estimated coefficient for the ruggedness interaction (scaled by the
standard error) when the observation is included and when it is
excluded. Following Belsley, Kuh, and Welsch (1980), we omit all
observations for which \(|\texttt{DFBETA}_i| > 2 / \sqrt{N}\), where
\(N\) is the number of observations, in our case 170.\footnote{Using
  other measures and rules for the omission of influential observations,
  such as dfits, Cook's distance or Welsch distance, provide very
  similar results} Results are presented in columns (3) and (4) of table
2.

In all four of the regressions with omitted observations, the ruggedness
coefficient remains negative and, except in column (1), statistically
significant. The ruggedness interaction remains positive and
statistically significant, confirming that there exists a differential
effect of ruggedness for Africa.

The reason why, in figure 2, a small number of observations appear as
particularly influential is because the ruggedness measure is skewed to
the left, leaving a small number of observations with large values. We
remedy this in two ways. First, we take the natural log of ruggedness
and use this in the estimating equations. This draws in the outlying
observations in the regression. The estimates of interest, reported in
columns (5) and (6), remain robust to this transformation. However,
looking at the natural log of ruggedness variable, one finds that the
measure is no longer left skewed, but it is now right skewed, with a
small number of influential observations taking on very small values.
Because of this, we pursue a second strategy where we perform a zero
skewness Box-Cox power transformation on the ruggedness variable to
obtain a measure with zero skewness. The relationships between income
and the zero-skewness ruggedness measure are shown in
Figure~\ref{fig-3}. It is evident that the income-ruggedness
relationships using zero-skewness measure do not feature influential,
outlying observations. In addition, a different relationship within
Africa and outside of Africa is still apparent in the scatter plots of
the data. Estimates using the zero-skewness measure are reported in
columns (7) and (8). The estimates confirm the impression given by the
figures. There is a positive and significant differential effect of
ruggedness within Africa.

\phantomsection\label{cell-fig-3}
\begin{figure}[H]

\centering{

\pandocbounded{\includegraphics[keepaspectratio]{index_files/mediabag/index_files/figure-pdf/fig-3-1.pdf}}

}

\caption{\label{fig-3}Results from Nunn and Puga (2012), where the
causal impact of terrain ruggedness on present-day GDP per capita
\textbf{reverses direction} when considering African vs.~non-African
countries!}

\end{figure}%

\textsubscript{Source:
\href{https://emilylmitchum.github.io/dsan5650-final-project/index.qmd.html}{Article
Notebook}}

\subsubsection{Is the effect of ruggedness also different for other
areas?}\label{is-the-effect-of-ruggedness-also-different-for-other-areas}

An obvious objection to the above results is that we have only allowed
the effect of ruggedness on economic outcomes to differ for African
countries. To check whether other continents of the world also
experience a differential impact of ruggedness, we add to our baseline
specification indicator variables for other continents as well as
interactions between ruggedness and these indicator variables. Estimates
are reported in table 3. The first column reports our baseline estimates
for comparison. Column (2) includes indicator variables for Europe,
North America, South America, and Oceania (the omitted continent is
Asia), and allows the effect of ruggedness to differ for each continent.
Overall, there is no evidence of differential effects for continents
other than Africa. The one exception is Europe, which is estimated to
have a positive differential ruggedness effect. However, this is driven
by just one country, Switzerland. When we omit Switzerland in column
(3), even the differential effect for Europe disappears.

A final possible source of concern is that perhaps the differential
effect of ruggedness for Africa is not really an African effect. Perhaps
it arises because the effect of ruggedness on income differs for areas
with some geographic characteristic that happens to be particularly
prevalent in Africa. For instance, it could be that in countries where a
large fraction of the territory experiences tropical climates, rugged
areas are cooler, dryer, or even less prone to tropical diseases. If
tropical climates are particularly prevalent in Africa (they
characterize 34.3\% of land in Africa compared with 19.2\% of the rest
of the world excluding Antarctica), perhaps the interaction between
ruggedness and the Africa indicator is proxying for an interaction
between ruggedness and tropical climates. Similarly, it could be that in
countries where a large fraction of the territory is covered by desert,
rugged areas are less arid. If deserts are particularly prevalent in
Africa (they cover 12.9\% of land in Africa compared with 3.4\% of the
rest of the world excluding Antarctica), perhaps the interaction between
ruggedness and the Africa indicator is proxying for an interaction
between ruggedness and desert.

We consider these possibilities in column (4) of table 3. There, we add
to our baseline estimation variables measuring the percentage of each
country with tropical climates and the percentage of each country
covered by desert (these can be seen as playing the same role as the
Africa indicator) as well as interactions between ruggedness and these
two variables (these can be seen as playing the same role as the
interaction between ruggedness and the Africa indicator). The
coefficients of interest measuring the direct effect of ruggedness and
the differential effect for Africa change little and remain
statistically significant.\footnote{One could also think that certain
  countries, because of inferior access to technology or poor
  governance, are worse equipped to mitigate the direct contemporaneous
  effect of ruggedness. However, note that this would work against
  estimating a positive differential effect of ruggedness in Africa
  since access to technology and governance are likely to be worse on
  average in Africa. A further concern is that the tropical climate
  measure is potentially endogenous to ruggedness, since some areas may
  not be classified as tropical if they are rugged. A preferable measure
  would quantify how tropical a climate would be if it were not rugged.
  Using a country's distance from the equator as a proxy for this
  measure yields very similar results.}

\subsection{Do slave exports account for Africa's differential
effect?}\label{sec-5}

Having established a robust differential impact of ruggedness within
Africa, we now examine whether Africa's slave trades can account for
this differential effect. We do this in two steps. First, we show that
there were fewer slaves exported from more rugged areas, even after we
control for initial population density and other potential geographic
determinants of slave exports. This involves estimating equation
(\textbf{?@eq-3}) from \textbf{?@sec-3}. Second, we estimate equation
(\ref{eq-7}) from \textbf{?@sec-3} to show that slave exports can fully
account for the differential effect of ruggedness in Africa.

We use data from Nunn (2008) on the number of slaves exported from each
country in Africa between 1400 and 1900 during Africa's four slave
trades. We normalize total slave exports by the land area of each
country. Because the variable is extremely skewed to the left and some
countries have zero slave exports, we take the natural logarithm of one
plus the measure, i.e.~\(\ln(1 + \text{slave exports} / \text{area})\).

Column (1) in table 4 reports the unconditional relationship between
slave exports and ruggedness for all 49 African countries included in
our previous estimations. It shows that there is a negative and
statistically significant relationship between ruggedness and slave
exports, and that ruggedness alone explains about one third of the
variation in slave exports. In the remaining columns of the table we
include additional variables to address several potential concerns
regarding the relationship between ruggedness and slave exports. First,
it is possible that the reason fewer slaves were taken from countries
with higher average uphill slope is that there were fewer people living
in more rugged areas, not because rugged terrain provided protection. To
allow for this possibility, in column (2) we control for population
density in 1500. We see that the negative relationship between
ruggedness and slave exports is robust, suggesting that the relationship
is not working through lower initial population density in rugged areas.
Next, to check whether ruggedness is proxying for other geographical
characteristics that could have affected slave exports, in column (3) we
control for the proportion of countries' land area that has tropical
climates, the proportion that is desert, and the average distance to the
coast. The negative relationship between slave exports and ruggedness
remains robust to these controls. The last set of controls that we
introduce are the distance from each country to the closest slave market
in each of the four slave trades (in thousands of kilometres). These
distances were shown by Nunn (2008) to be important determinants of
slave exports. Again, the ruggedness coefficient remains negative and
significant when these distances are controlled for. In the final column
of the table we include all control variables simultaneously. The
significant negative relationship between slave exports and ruggedness
remains. Overall, the results of table 4 confirm the historical accounts
suggesting that irregular terrain was an effective means of avoiding the
slave trades.

Having shown that ruggedness has a differential effect in Africa and
that ruggedness deterred slave exports, we now check to what extent
slave exports can account for Africa's differential effect. We do this
by estimating equation (\ref{eq-7}) from \textbf{?@sec-3}, with results
reported in table 5. Column (1) reproduces, for comparison, our
estimates of equation (5), which does not control for slave exports. In
columns (2), we estimate equation (7), which is identical to equation
(5), except that slave exports are also included in the estimating
equation. Columns (3) and (4) report the same estimations as (1) and (2)
respectively with the same controls used in Section~\ref{sec-4}. With or
without the full set of controls, when slave exports are controlled for,
the estimated differential effect of ruggedness within Africa
disappears. The estimated coefficient is essentially zero, and is no
longer statistically significant. This confirms hypothesis 5, and
provides support for the explanation that the differential effect of
ruggedness arises because of the slave trades. The estimates also
support hypotheses 4 and 6. Slave exports have had a negative and
statistically significant effect on economic outcomes, and once slave
exports are accounted for, the estimated direct effect of ruggedness on
income is negative and statistically significant.

\subsection{The distribution of population relative to rugged
areas}\label{the-distribution-of-population-relative-to-rugged-areas}

We have shown that the data support our hypothesis that, within Africa,
rugged terrain has had an indirect benefit on current economic outcomes
because inhabitants of rugged areas were better able to protect
themselves from the slave trades and avoid their long-run negative
consequences. In this section, we also show that this benefit may have
come at a cost, since it has meant that African populations are now more
concentrated in rugged areas relative to populations in other parts of
the world.

\subsubsection{Migration towards rugged areas in response to the slave
trades}\label{migration-towards-rugged-areas-in-response-to-the-slave-trades}

In parts of Africa threatened by the slave trades, populations in more
rugged areas tended to increase relative to populations in less rugged
areas. This is because in rugged areas people were better able to defend
against the depopulating effects of the slave trades, and also because
population tended to migrate to more rugged areas escaping from slave
raids (see Cordell (2004)). If these relocations were permanent, then
populations today will be more concentrated in rugged areas, which we
have seen adversely affects income.

There are a number of reasons why past population movements may persist.
If there are benefits to agglomeration, then relocation requires a
coordinated effort, which may be difficult. As well, many forms of
capital, such as wells, buildings, and streets, are very difficult or
even impossible to move. Finally, culture and religious beliefs that
incorporate the rugged landscape may impede relocation. For example, the
Dogon, who inhabit the cliffs of the Badiagara escarpment in Mali, place
great spiritual significance on the caves located on the cliffs. They
believe that spirits inhabit the cliffs, and they also bury their dead
there (Olson, 1996).

We now test whether there is evidence that these historically documented
population movements persist until today. We do this by constructing a
measure of whether people in a country tend to live in the more or less
rugged parts of the country. Consider again \textbf{?@fig-1}, which
shows the construction of average ruggedness for each 30 arc-second by
30 arc-second cell within a country. Combining these ruggedness data
with data from Oak Ridge National Laboratory (2001) on the total ambient
population in each cell of the same grid, we calculate the correlation
coefficient between population density and terrain ruggedness in each
country using 30 by 30 arc-second cells as units (see the appendix for
further details). Not surprisingly, we find that for most countries (138
of 170), the correlation coefficient is negative, indicating that
populations tend to concentrate in the less rugged parts of the country.

Using our population-ruggedness correlation measure, we first examine
whether there is evidence that within Africa the combination of the
threat of slave raids with the availability of areas with varying
ruggedness created past population movements that have persisted until
today. In particular, we estimate the following equation:

\begin{equation}\phantomsection\label{eq-8}{
\text{Corr}(\text{Pop. dens}, \text{Ruggedness}) = \eta_0 + \eta_1 \sigma_i^r + \eta_2 x_i + \eta_3 x_i \sigma_i^r + \varepsilon_i,
}\end{equation}

where \(\text{Corr}(\text{Pop. dens}, \text{Ruggedness})\) is our
ruggedness-population density correlation measure, \(\sigma^r_i\) is the
standard deviation of ruggedness within each country, and as before
\(x_i\) denotes slave exports.

This specification captures the effect that the threat of capture during
the slave trades had on population movements. We take our measure of
slave exports \(x_i\) as a proxy for how exposed a country was to the
slave trades. If a country is threatened by the slave trades, then this
may cause migration to more rugged areas. However, this will only be
feasible if more rugged areas are available to migrate to. Therefore,
the movement will depend not only on the threat from the slave trades
\(x_i\), but also on the variation in terrain ruggedness within the
country \(\sigma^r_i\). The importance of the interaction between the
threat from the slave trades and varying ruggedness for population
patterns is captured by the coefficient \(\eta_3\) that multiplies
\(x_i \sigma^r_i\) in the estimating equation. Because slave exports are
expected to have a greater effect on migration when there are more
differences in terrain ruggedness within a country, we expect
\(\eta_3 > 0\). At the extreme, if the whole country is equally rugged
(\(\sigma^r_i = 0\)), then there is no benefit to moving to another part
of the country in response to the threat of capture. In this case, the
relationship between population concentration and slave exports in
(\ref{eq-8}) is given by \(\eta_2\) only, and therefore we expect
\(\eta_2 = 0\).

The estimates of Equation~\ref{eq-8} are reported in column (1) of table
6. As expected, the estimated coefficient for slave exports alone
\(\eta_2\) is not significantly different from zero, but the estimated
coefficient for slave exports interacted with the standard deviation of
ruggedness \(\eta_3\) is positive and statistically significant. In
column (2) we also control for our set of country geographic
characteristics. The results remain robust to this. Overall, the
estimates suggest that when populations within African countries were
threatened by the slave trades and differences in ruggedness existed,
then less rugged areas lost population relative to more rugged areas.
This concentration of the population in the more rugged parts of the
country persists today.

\subsubsection{Controlling for the distribution of population relative
to rugged
areas}\label{controlling-for-the-distribution-of-population-relative-to-rugged-areas}

This evidence that the slave trades caused persistent movements to more
rugged areas suggests an additional effect of ruggedness within Africa.
Although ruggedness provided protection against the slave trades, this
protection came at a cost. Populations became permanently more
concentrated in rugged areas. An implication of this is that our
estimate of the indirect effect of ruggedness so far (the estimated
coefficient for the interaction between ruggedness and the African
indicator variable) has captured both the long-run benefit from
ruggedness as a deterrent of slave exports, as well as the long-run cost
from having population persistently more concentrated in rugged areas in
response to the threat of the slave trades. We now try to isolate the
benefit from ruggedness as a deterrent of slave exports, by controlling
for concentration of population in more rugged areas. We do this through
two complementary strategies.

First, we include the ruggedness-population correlation measure in
estimating \textbf{?@eq-5}, while allowing the effect this variable to
differ within Africa and outside of Africa. The results are reported in
column (3) and (4) of table 7 (where column (4) incorporates the same
controls used in Section~\ref{sec-4}). Columns (1) and (2) reproduce our
baseline estimates of Table~\ref{tbl-1} for comparison. Consistent with
our expectations, when one controls for the current population
concentration, the estimated indirect effect of ruggedness (i.e., the
coefficient for \(\text{ruggedness}\cdot I^\text{Africa}\)) increases.
In other words, the slave trades have not only created a positive
indirect effect of ruggedness on income in Africa, but by encouraging
population concentrations in rugged areas, they have also amplified the
negative direct impact of ruggedness in Africa. However, the differences
between estimates in columns (1) and (2) and between the estimates of
columns (3) and (4) are not large, which suggests that the indirect
benefits of ruggedness by deterring slave exports in Africa are much
greater than the costs that arise from past migrations to more rugged
areas.

Second, rather than using a separate measure to control for a country's
concentration of its populations in rugged areas, an alternative
strategy is to construct a ruggedness measure that weights mean uphill
slope by population rather than by land area.\footnote{As detailed in
  \textbf{?@sec-2}, to construct this population-weighted ruggedness
  measure, we start from the same mean uphill slope for each 30 by 30
  arc-second cell as in our standard ruggedness measure but, in
  averaging this for each country, we weight ruggedness in each 30 by 30
  arc-second cell by the share of the country's population located in
  that 30 by 30 arc-second cell in 2000.} This allows the
contemporaneous impact of ruggedness on income to work locally in
proportion to population density, so that ruggedness matters more in
densely populated areas and not at all in uninhabited areas. The
estimates using this alternative ruggedness measure are reported in
columns (5) and (6).\footnote{The specifications in columns (5) and (6)
  are identical to those in columns (1) and (2) respectively, with our
  standard ruggedness measure replaced by population-weighted
  ruggedness. The measure of ruggedness interacted with the African
  indicator variable in columns (5) and (6) is still our standard
  ruggedness variable, since there is no reason why the strength of the
  historical channel through which ruggedness affects income should be
  proportional to current population densities. Nevertheless,
  interacting population-weighted ruggedness instead of area-weighted
  ruggedness with the African indicator variable does not affect our
  qualitative results. As noted earlier, we have chosen to use
  area-weighted ruggedness up until this last section because it is
  truly exogenous (whereas the distribution of population is potentially
  related to various economic, cultural, and institutional factors) and
  in order to use a common measure to estimate the direct
  contemporaneous and indirect historical effects of ruggedness.} The
direct effect remains negative and significant and the differential
effect in Africa remains positive, larger in magnitude than the direct
effect, and statistically significant. Once again, comparison with
columns (1) and (2) shows that, once we account for the current
population concentration across areas with different ruggedness, the
estimated indirect effect of ruggedness increases.

\subsection{Conclusions}\label{conclusions}

This paper contributes to the ongoing controversy about whether
geography matters mainly because of its direct contemporaneous impact on
economic outcomes or because of its interaction with historical events.
The issue is not just a matter of intellectual curiosity. It is very
important for development policies, and has been the source of much
recent debate. Researchers such as Sachs (2005) and others, who believe
that geography matters primarily through a direct contemporaneous
channel, argue that foreign aid and investment are needed to overcome
the adverse geographic environments of poor countries. According to this
view, with enough foreign aid the adverse effects of geography can be
mitigated and alleviated, allowing the world's poorest countries to
develop. Arguing against the effectiveness of foreign aid are those such
as William Easterly (2006a,b, 2007), who believe that the importance of
geography lies in its influence on past events, which have shaped the
evolution of societies.\footnote{See in particular Easterly and Levine
  (2003).} It is felt that the reason for underdevelopment has more to
do with poor domestic governance, dysfunctional institutions, or poor
economic policies, all of which have deep historical roots. According to
this view, aid rather than being a panacea for underdevelopment, is
largely ineffective and can even exacerbate poor governance, hampering
economic development.\footnote{See Djankov, Montalvo, and Reynal-Querol
  (2008) for recent evidence on this ``curse of aid''.}

By focusing on a dimension of geography, terrain ruggedness, which
varies throughout the world and on a historical event, the slaves trades
over the period 1400-1900, which is geographically confined to Africa,
we are able to separately identify the direct contemporaneous channel
and the indirect historic channel. We find a direct negative effect of
ruggedness on income, which is consistent with irregular terrain making
agriculture, building, and transportation more costly. We also find that
rugged terrain had an additional effect in Africa during the 15th to
19th centuries: it afforded protection to those being raided during
Africa's slave trades. By allowing areas to escape from the detrimental
effects that the slave trades had on subsequent economic development,
ruggedness also creates long-run benefits in Africa through an indirect
historic channel. This differential effect of ruggedness is confined to
Africa, is not due to geographic features that are particularly
prominent in this continent, and can be fully accounted for by the slave
trades.

On the whole, the results point to the importance of both a direct and
an indirect channel through which ruggedness affects income. In terms of
relative magnitudes, within Africa, the positive indirect effect of
ruggedness swamps the direct negative effect. In this case, geography
matters more through its indirect effect on past events, than it does
through it current direct effect on income. However, these results are
specific to terrain ruggedness and do not imply that the indirect
historic channel dominates the direct contemporaneous channel in other
contexts.

Finally, our paper also suggests that the direct and indirect channels
can interact in interesting ways. The combination of the threat of slave
raids with the availability of areas with varying ruggedness has
resulted in a concentration of population in particularly rugged areas
that has persisted until today. In doing so, in addition to creating an
indirect benefit of ruggedness in Africa, the slave trades have also
amplified the negative direct and contemporaneous impact of ruggedness
in Africa.

\subsection{References}\label{references}

\phantomsection\label{refs}
\begin{CSLReferences}{1}{0}
\bibitem[\citeproctext]{ref-belsley_regression_1980}
Belsley, David A, Edwin Kuh, and Roy E Welsch. 1980. \emph{Regression
Diagnostics: {Identifying} Influential Data and Sources of
Collinearity}. New York: {John Wiley and Sons}.

\bibitem[\citeproctext]{ref-cordell_myth_2004}
Cordell, Dennis D. 2004.
{``\href{https://www.cambridge.org/core/books/fighting-the-slave-trade/myth-of-inevitability-and-invincibility-resistance-to-slavers-and-the-slave-trade-in-central-africa-18501910/E2E9819E2BBD9A958F4BB7CDC505764C}{The
{Myth} of {Inevitability} and {Invincibility}: {Resistance} to {Slavers}
and the {Slave Trade} in {Central Africa}, 1850--1910}.''} In
\emph{Fighting the {Slave Trade}: {West African Strategies}}, edited by
Sylviane A. Diouf, 31--49. Boydell \& Brewer.

\bibitem[\citeproctext]{ref-djankov_curse_2008}
Djankov, Simeon, Jose G. Montalvo, and Marta Reynal-Querol. 2008.
{``\href{https://doi.org/10.1007/s10887-008-9032-8}{The Curse of
Aid}.''} \emph{Journal of Economic Growth} 13 (3): 169--94.

\bibitem[\citeproctext]{ref-easterly_tropics_2003}
Easterly, William, and Ross Levine. 2003.
{``\href{https://doi.org/10.1016/S0304-3932(02)00200-3}{Tropics, Germs,
and Crops: How Endowments Influence Economic Development}.''}
\emph{Journal of Monetary Economics} 50 (1): 3--39.

\bibitem[\citeproctext]{ref-nunn_longterm_2008}
Nunn, Nathan. 2008.
{``\href{https://doi.org/10.1162/qjec.2008.123.1.139}{The {Long-Term
Effects} of {Africa}'s {Slave Trades}}.''} \emph{Quarterly Journal of
Economics} 123 (1): 139--76.

\bibitem[\citeproctext]{ref-nunn_ruggedness_2012}
Nunn, Nathan, and Diego Puga. 2012.
{``\href{https://www.jstor.org/stable/41349158}{Ruggedness: {The
Blessing} of {Bad Geography} in {Africa}}.''} \emph{The Review of
Economics and Statistics} 94 (1): 20--36.

\bibitem[\citeproctext]{ref-sachs_end_2005}
Sachs, Jeffrey. 2005.
\emph{\href{https://books.google.com?id=uzhveDZnv_gC}{The {End} of
{Poverty}: {How We Can Make} It {Happen} in {Our Lifetime}}}. Penguin
Books Limited.

\end{CSLReferences}




\end{document}
